% GENERATED FILE. Edit canonical lessons, then run npm run generate:books.
% canonical-source-hash: fnv1a64:d3455fe80f2237d4
% canonical-lessons: TE-C217-kalisi, TE-C217-tappa, TE-C217-sampradayam, TE-C217-acaram, TE-C217-alavatu

\chapter{Together, Except, Tradition, Custom, Habit}
\label{ch:te-together-except-tradition-custom-habit}
\hlchaptermodality{217}

\begin{chapteropening}
\textbf{By the end of this chapter:} \emph{I can say together, except, a tradition, a custom and a habit.}

Complete the last lesson of chapter 217: I can say together, except, a tradition, a custom and a habit.
\end{chapteropening}

\section[\texorpdfstring{\te{కలిసి} (kalisi)}{కలిసి (kalisi)}]{\te{కలిసి} (kalisi) — together}

\label{lesson:TE-C217-kalisi}



\begin{quote}
Before the new one: say the Telugu for a bakery, then the Telugu for a hair salon.
\end{quote}

\subsection*{You'll want to know: \te{కలిసి}}
\textbf{\te{కలిసి}} — \emph{kalisi} — \textquotedblleft{}together\textquotedblright{}.

\begin{grammarlens}[title={What you've built}]
One more word for this chapter.
\end{grammarlens}

\subsection*{Your turn}
\begin{itemize}
\raggedright
  \item \emph{Say it:} \emph{kalisi}
  \item \emph{Say it:} \emph{kalisi}, once more
  \item \emph{Say it:} say \emph{kalisi}
  \item \emph{Recall:} say the Telugu for a bakery, then the Telugu for a hair salon, then say \emph{kalisi} again
\end{itemize}

\begin{tcolorbox}[breakable,colback=teal!4,colframe=teal!35!black,title={Before you move on}]
What does \te{కలిసి} mean? (\textquotedblleft{}Together\textquotedblright{}.) Say it once more.
\end{tcolorbox}
\section[\texorpdfstring{\te{తప్ప} (tappa)}{తప్ప (tappa)}]{\te{తప్ప} (tappa) — except}

\label{lesson:TE-C217-tappa}



\begin{quote}
Before the new one: say the Telugu for a hair salon, then the Telugu for together.
\end{quote}

\subsection*{You'll want to know: \te{తప్ప}}
\textbf{\te{తప్ప}} — \emph{tappa} — \textquotedblleft{}except\textquotedblright{}.

\begin{grammarlens}[title={What you've built}]
One more word for this chapter.
\end{grammarlens}

\subsection*{Your turn}
\begin{itemize}
\raggedright
  \item \emph{Say it:} \emph{tappa}
  \item \emph{Say it:} \emph{tappa}, once more
  \item \emph{Say it:} say \emph{tappa}
  \item \emph{Recall:} say the Telugu for a hair salon, then the Telugu for together, then say \emph{tappa} again
\end{itemize}

\begin{tcolorbox}[breakable,colback=teal!4,colframe=teal!35!black,title={Before you move on}]
What does \te{తప్ప} mean? (\textquotedblleft{}Except\textquotedblright{}.) Say it once more.
\end{tcolorbox}
\section[\texorpdfstring{\te{సంప్రదాయం} (sampradāyaṁ)}{సంప్రదాయం (sampradāyaṁ)}]{\te{సంప్రదాయం} (sampradāyaṁ) — a tradition}

\label{lesson:TE-C217-sampradayam}



\begin{quote}
Before the new one: say the Telugu for together, then the Telugu for except.
\end{quote}

\subsection*{You'll want to know: \te{సంప్రదాయం}}
\textbf{\te{సంప్రదాయం}} — \emph{sampradāyaṁ} — \textquotedblleft{}a tradition\textquotedblright{}.

\begin{grammarlens}[title={What you've built}]
One more word for this chapter.
\end{grammarlens}

\subsection*{Your turn}
\begin{itemize}
\raggedright
  \item \emph{Say it:} \emph{sampradāyaṁ}
  \item \emph{Say it:} \emph{sampradāyaṁ}, once more
  \item \emph{Say it:} say \emph{sampradāyaṁ}
  \item \emph{Recall:} say the Telugu for together, then the Telugu for except, then say \emph{sampradāyaṁ} again
\end{itemize}

\begin{tcolorbox}[breakable,colback=teal!4,colframe=teal!35!black,title={Before you move on}]
What does \te{సంప్రదాయం} mean? (\textquotedblleft{}A tradition\textquotedblright{}.) Say it once more.
\end{tcolorbox}
\section[\texorpdfstring{\te{ఆచారం} (ācāraṁ)}{ఆచారం (ācāraṁ)}]{\te{ఆచారం} (ācāraṁ) — a custom, a practice}

\label{lesson:TE-C217-acaram}



\begin{quote}
Before the new one: say the Telugu for except, then the Telugu for a tradition.
\end{quote}

\subsection*{You'll want to know: \te{ఆచారం}}
\textbf{\te{ఆచారం}} — \emph{ācāraṁ} — \textquotedblleft{}a custom, a practice\textquotedblright{}.

\begin{grammarlens}[title={What you've built}]
One more word for this chapter.
\end{grammarlens}

\subsection*{Your turn}
\begin{itemize}
\raggedright
  \item \emph{Say it:} \emph{ācāraṁ}
  \item \emph{Say it:} \emph{ācāraṁ}, once more
  \item \emph{Say it:} say \emph{ācāraṁ}
  \item \emph{Recall:} say the Telugu for except, then the Telugu for a tradition, then say \emph{ācāraṁ} again
\end{itemize}

\begin{tcolorbox}[breakable,colback=teal!4,colframe=teal!35!black,title={Before you move on}]
What does \te{ఆచారం} mean? (\textquotedblleft{}A custom, a practice\textquotedblright{}.) Say it once more.
\end{tcolorbox}
\section[\texorpdfstring{\te{అలవాటు} (alavāṭu)}{అలవాటు (alavāṭu)}]{\te{అలవాటు} (alavāṭu) — a habit}

\label{lesson:TE-C217-alavatu}



\begin{quote}
Before the new one: say the Telugu for a tradition, then the Telugu for a custom.
\end{quote}

\subsection*{You'll want to know: \te{అలవాటు}}
\textbf{\te{అలవాటు}} — \emph{alavāṭu} — \textquotedblleft{}a habit\textquotedblright{}.

\begin{grammarlens}[title={What you've built}]
One more word for this chapter.
\end{grammarlens}

\subsection*{Your turn}
\begin{itemize}
\raggedright
  \item \emph{Say it:} \emph{alavāṭu}
  \item \emph{Say it:} \emph{alavāṭu}, once more
  \item \emph{Say it:} say \emph{alavāṭu}
  \item \emph{Recall:} say the Telugu for a tradition, then the Telugu for a custom, then say \emph{alavāṭu} again
\end{itemize}

\begin{tcolorbox}[breakable,colback=teal!4,colframe=teal!35!black,title={Before you move on}]
What does \te{అలవాటు} mean? (\textquotedblleft{}A habit\textquotedblright{}.) Say it once more.
\end{tcolorbox}
